\documentclass[10pt, twocolumn]{article}
\usepackage{geometry}
\geometry{a4paper, margin=1in}
\usepackage{graphicx}
\usepackage{amsmath}
\usepackage{amssymb}
\usepackage{algorithm}
\usepackage{algpseudocode}
\usepackage{hyperref}
\usepackage{booktabs}
\usepackage{caption}
\usepackage{subcaption}

\title{\textbf{Robust Doorway Counting and Navigation via Integration of POMDP Models and PID Controller States}}
\author{Autonomy Lab Research Team \\ \textit{Department of Computer Science and Engineering}}
\date{\today}

\begin{document}

\maketitle

\begin{abstract}
Autonomous navigation in corridor environments presents a seemingly simple yet noisy state estimation challenge: distinguishing between feature-poor walls and open doorways while maintaining stable motion. This paper presents a novel hybrid architecture that fuses a continuous Proportional-Integral-Derivative (PID) controller with a discrete Partially Observable Markov Decision Process (POMDP). Unlike traditional approaches that rely solely on raw range sensor data, our system incorporates the internal error dynamics of the PID controller itself as a key feature in the observation space. We model the corridor as a linear chain of hidden states and employ a Naive Bayes observation model to infer the robot's location. Real-time inference is performed on an iRobot Create 3 platform, where a decision policy based on expected reward maximization successfully triggers a return-to-home maneuver after counting a specific sequence of doorways. Experimental results demonstrate that fusing controller dynamics with odometry and proximity data significantly improves state disambiguation compared to proximity sensors alone.
\end{abstract}

\section{Introduction}
The problem of "corridor navigation" is a canonical task in mobile robotics, often serving as a proxy for broader challenges in simultaneous localization and mapping (SLAM) and semantic place recognition. While modern LiDAR-based SLAM systems can map environments with high precision, they are often computationally intensive and expensive. For low-cost platforms restricted to sparse sensors (e.g., simple IR proximity or sonar), the challenge lies in reliably identifying landmarks (doorways) amidst sensor noise and drift.

A common failure mode in reactive wall-following is the inability to distinguish between a true doorway (which requires no corrective action or a specific "ignore" behavior) and a sensor glitch or wall irregularity. Furthermore, strictly reactive systems lack a notion of "global progress," making tasks like "pass three doors then return" difficult to execute robustly.

In this work, we propose a probabilistic framework that couples low-level control with high-level reasoning. Our contribution is twofold:
\begin{enumerate}
    \item A \textbf{POMDP-based state estimator} that models the corridor as a topological chain, allowing the robot to maintain a belief distribution over its progress even when individual sensor readings are ambiguous.
    \item A \textbf{Novel Feature Space} that treats the internal states of the PID wall-following controller (Proportional, Integral, and Derivative error terms) as "virtual sensors." We hypothesize that the controller's reaction to a doorway (e.g., a rapid derivative spike followed by integral windup) provides a distinct signature that aids classification.
\end{enumerate}

\section{System Overview}
\subsection{Hardware Platform}
The system is implemented on the \textbf{iRobot Create 3}, a differential drive mobile robot. The sensor suite utilized for this study includes:
\begin{itemize}
    \item \textbf{IR Proximity Sensors:} specifically the side-facing sensors, calibrated to measure distance to the wall in centimeters.
    \item \textbf{Wheel Encoders:} providing odometry for estimating distance traveled (\(\Delta d\)) between time steps.
    \item \textbf{Bump Sensors:} providing binary contact information for collision recovery.
\end{itemize}

\subsection{Software Architecture}
The control stack is written in Python using `asyncio` for concurrent execution of the control loop (10Hz) and the belief update mechanism. The architecture is hierarchical:
\begin{itemize}
    \item \textbf{Level 1 (Reactive):} A PID controller maintains a setpoint distance of 6.0 cm from the right wall.
    \item \textbf{Level 2 (Probabilistic):} A POMDP belief filter updates the probability of the robot's location every 10 cm of travel.
    \item \textbf{Level 3 (Executive):} A policy monitors the Expected Reward of the current belief to trigger mission events (e.g., turning around).
\end{itemize}

\section{Methodology}

\subsection{The PID Controller as a Sensor}
Standard wall-following uses a PID controller to minimize the error \(e(t) = d_{setpoint} - d_{measured}\).
\begin{equation}
 u(t) = K_p e(t) + K_i \int_0^t e(\tau) d\tau + K_d \frac{de(t)}{dt}
\end{equation}
Typically, \(u(t)\) is sent to the motors and \(e(t)\) is discarded. However, we observe that \(e(t)\), \(\int e\), and \(\frac{de}{dt}\) behave predictably in different environmental contexts:
\begin{itemize}
    \item \textbf{Wall:} \(e(t) \approx 0\), derivative is low.
    \item \textbf{Door Start:} Sharp increase in \(e(t)\) and positive derivative spike.
    \item \textbf{Door:} Large sustained \(e(t)\), integral term accumulates (windup).
\end{itemize}
We discretize these internal terms into 10 bins and feed them as observations into the Bayesian network, effectively allowing the robot to "feel" the shape of the wall through its own control efforts.

\subsection{POMDP Formulation}
We model the environment as a Partially Observable Markov Decision Process (POMDP) defined by the tuple \((S, A, T, R, \Omega, O)\).

\subsubsection{State Space (\(S\)}
The corridor is modeled as a linear sequence of discrete states. To capture the transient dynamics of passing a door, we subdivide the "Door" event into three phases. For a mission requiring the passage of 3 doors, the state space is:
\begin{align*}
S = \{ &\text{Wall}_0, \\
       &\text{DoorStart}_1, \text{Door}_1, \text{DoorPassed}_1, \\
       &\text{Wall}_1, \\
       &\dots \\
       &\text{Wall}_{End} \}
\end{align*}
This fine-grained discretization allows the model to anticipate the end of a door and the resumption of a wall.

\subsubsection{Transition Model (\(T\)}
The transition function \(T(s' | s, a)\) describes the probability of moving from state \(s\) to \(s'\). Since the robot's primary action is "move forward," we simplify this to a motion model based on distance traveled. We define a "stay probability" \(P_{stay} = 0.8\) and a "progression probability" \(P_{step} = 0.2\).
\begin{equation}
P(s_{t+1} | s_t) = \begin{bmatrix}
0.8 & 0.2 & 0 & \dots \\
0 & 0.8 & 0.2 & \dots \\
\vdots & \vdots & \ddots & \vdots \\
0 & 0 & \dots & 1.0
\end{bmatrix}
\end{equation}
The final state, \(\text{Wall}_{End}\), is an absorbing state.

\subsubsection{Observation Model (\(O\)}
The observation vector \(z_t\) at time \(t\) is a collection of \(N\) discrete features:
\begin{equation}
z_t = \{ IR_{1:9}, PID\_P_{1:9}, PID\_I_{1:9}, PID\_D_{1:9}, ODO_{1:9} \}
\end{equation}
We utilize a history of length 9 for each feature to capture temporal trends. To make inference tractable, we employ a Naive Bayes assumption, treating features as conditionally independent given the state:
\begin{equation}
P(z_t | s_t) = \prod_{f \in z_t} P(f | s_t)
\end{equation}
The individual probabilities \(P(f | s_t)\) are stored in Conditional Probability Tables (CPTs) learned from training data.

\subsubsection{Belief Update}
The belief state \(b_t(s)\) is the posterior probability distribution over \(S\). It is updated recursively:
\begin{enumerate}
    \item \textbf{Prediction Step (Motion):}
    \begin{equation}
    \bar{b}_t(s') = \sum_{s \in S} T(s' | s) b_{t-1}(s)
    \end{equation}
    \item \textbf{Correction Step (Observation):}
    \begin{equation}
    b_t(s) = \eta P(z_t | s) \bar{b}_t(s)
    \end{equation}
\end{enumerate}
where \(\eta = 1 / \sum_{s} P(z_t | s) \bar{b}_t(s)\) is the normalization constant.

\subsection{Decision Policy}
Instead of acting on the most likely state (\(\arg\max_s b_t(s)\)), which can oscillate, we calculate the \textbf{Expected Reward} to determine mission completion. We define a sparse reward function \(R(s)\) where \(R(s)=1\) for states indicative of mission success (e.g., high probability of having passed 3 doors) and 0 otherwise.
\begin{equation}
E[R]_t = \sum_{s \in S} b_t(s) R(s)
\end{equation}
The robot executes the "Return Home" maneuver only when \(E[R]_t > \tau\), with \(\tau = 0.8\).

\section{Experiments}
\subsection{Training}
Training data was collected by manually driving the robot along the corridor wall. A human operator annotated the ground truth location (e.g., pressing 'd' for Door, 'w' for Wall) in real-time. This resulted in labeled time-series data for all sensors. The CPTs were generated by discretizing the continuous sensor values into 12 bins and computing the frequency counts for each state-observation pair.

\subsection{Validation}
The system was deployed in a physical test environment containing three doorways spaced irregularly. The robot was tasked with autonomously following the wall, counting three doors, and turning around.

\section{Results and Discussion}
The integration of PID states proved crucial. In segments where the wall material changed reflectivity (confusing the IR sensors), the PID derivative term remained stable (as the physical distance hadn't changed), preventing false "Door" detections.

The belief distribution evolved as expected:
\begin{enumerate}
    \item \textbf{Initial Uncertainty:} At startup, belief was concentrated on \(Wall_0\).
    \item \textbf{Door Event:} Upon encountering the first gap, probability mass shifted rapidly to \(DoorStart_1\) and then \(Door_1\).
    \item \textbf{Transition:} As the robot re-acquired the wall, the \(DoorPassed_1\) state peaked, confirming the count.
\end{enumerate}
The expected reward metric provided a smooth, monotonic signal for the final decision, triggering the turn-around maneuver precisely after the third door was cleared, avoiding the "jitter" often seen in threshold-based logic.

\section{Conclusion}
We have presented a robust doorway counting framework that leverages the synergy between low-level control dynamics and high-level probabilistic reasoning. By treating the PID controller as a sensor and maintaining a full belief distribution via a POMDP, the system achieves reliable navigation in the presence of significant sensor noise. Future work will explore learning transition probabilities online to adapt to varying corridor lengths and crowd dynamics.

\end{document}